\section{Proof of \texorpdfstring{\cref{thm:routing}}{Staged step learning}: a schedule learns the routing}
\label{app:routing}

\begin{restated}{\Cref{thm:routing}}
For counting modulo~$K$ with the gate head, started from the small structured initialization below, staged block-coordinate
gradient flow on $\Lstep$ enters $\mathcal G_\zeta$ by a time $T\le\mathrm{poly}(n,K,d_{\mathrm h},\initscale^{-1},(p(1-p))^{-1})$,
stays in it at every later time, and drives $\Lstep$ to zero, for every $0<\initscale\le\initscale_\star$ with
$1/\initscale_\star\le C\sqrt w$.
\end{restated}

\paragraph{Roadmap.} The structured initialization reduces the trainable network to two $K\times K$ readout blocks and one
scalar gate feature $M_t$. Stage A fits the count-only prior; Stage B gives the gate readout a small coefficient on the move
$\bm e_{s+1}-\bm e_s$; Stage C trains the relative bias so that attention concentrates on the displayed count and current letter;
Stage D freezes the routing and trains the readout to zero loss. We prove each stage in turn and then combine their stopping times.
Symbols are local to this section.

\subsection{Reduced gate model and structured seed}
We count modulo $K$ over $\{q,r\}$, with a fixed query of probability $p\in(0,1)$ and a uniform start.

\emph{Task notation.} Write $\bm x_t=\bm e_{s_{t-1}}$ for the displayed count and $m_t=(-1)^{b_t}$ for the letter sign. Since $s_0$ is
uniform and independent of the letters, $s_{t-1}$ is uniform and independent of all $n$ letters. Let $R$ be the cyclic shift,
$R\bm e_s=\bm e_{s+1}$, so that $(R-I)\bm e_s=\bm e_{s+1}-\bm e_s$ is the move of the count when $b_t=1$, and recall the state-only predictor
$\bm q_s=(1-p)\bm e_s+p\,\bm e_{s+1}$. For $\bm v\in\R^K$ write $\operatorname{osc}(\bm v)=\max_k v_k-\min_k v_k$. The margin of a
prediction with logits $\bm z_t$ is $\gamma_t=z_{t,s_t}-\max_{k\ne s_t}z_{t,k}$, and $\sigma(x)=(1+e^{-x})^{-1}$ is the logistic sigmoid.

\paragraph{Small structured initialization.} Fix orthonormal residual directions $\hat e_{0},\hat e_\ell$, a count register
$\hat e^{\mathrm c}_0,\ldots,\hat e^{\mathrm c}_{K-1}$ and a gate register $\hat e^{\mathrm g}_0,\ldots,\hat e^{\mathrm g}_{K-1}$, and a unit
vector $\hat e_{\mathrm h}$ in the head space. With independent random signs $A_Q,A_K,A_V\in\{\pm1\}$, set
\begin{gather*}
 E(s)=\initscale\,\hat e^{\mathrm c}_s,\qquad
 E(q)=\initscale(\hat e_{0}-\hat e_\ell),\qquad E(r)=\initscale(\hat e_{0}+\hat e_\ell),\\
 W_Q\hat e^{\mathrm c}_s=\initscale A_Q\hat e_{\mathrm h},\qquad W_K\hat e_\ell=\initscale A_K\hat e_{\mathrm h},\qquad
 W_V\hat e^{\mathrm c}_s=\initscale A_V\hat e^{\mathrm g}_s\qquad(s\in\mathbb Z_K)
\end{gather*}
for the gate head (head $1$ of block $2$), with these maps zero on all other directions. Position embeddings, the other heads, the MLP
weights and the readout are zero, and the MLP biases are $\pm\initscale$. This $\theta_0$ satisfies \cref{ass:init}: every block
has operator norm at most $\sqrt K\,\initscale$, $|\tilde b_m|=1$, and the count embeddings have singular values $\initscale$.

The seed gives gate strength $a$ and value scale $g$,
\[
 a=\frac{\initscale^4A_QA_K}{\sqrt{d_{\mathrm h}}},\qquad g=\initscale^2A_V,
 \qquad |a|=\frac{\initscale^4}{\sqrt{d_{\mathrm h}}},\quad |g|=\initscale^2.
\]
At the position $i_t$ predicting $s_t$ ($t\ge2$), the gate head uses the segment mask: it attends to the displayed count at $i_t$
and the $n$ letters at $\ell_1,\ldots,\ell_n$. The current letter has relative offset $i_t-\ell_t=n$.
The key reads the letter and the value of the count position carries the displayed count, so the head multiplies a count by a
letter-dependent attention weight; the relative biases determine which sites are selected.
Let $U$ and $W$ be the $K\times K$ readout blocks on the count and gate registers; the rest of the readout stays zero.
The state logits have the following exact form.

\begin{lemma}[Exact features]
\label{lem:R-features}
Throughout the schedule, for $t\ge2$,
\[
 \bm z_t=\initscale U\bm x_t+g\,(1-M_t)\,W\bm x_t,
 \qquad
 M_t=\frac{\sum_{j=1}^ne^{am_j+\beta_{i_t-\ell_j}}}{e^{\beta_0}+\sum_{j=1}^ne^{am_j+\beta_{i_t-\ell_j}}}.
\]
In particular $M_t$ depends on the letters and $\beta$ but not on the counts.
\end{lemma}

\begin{proof}
Block $1$ and both MLPs write nothing. The query at $i_t$ is $\initscale^2A_Q\hat e_{\mathrm h}$ for every displayed count, the key of
letter $j$ is $\initscale^2A_Km_j\hat e_{\mathrm h}$ and that of $i_t$ is $0$, so letter $j$ has score $am_j+\beta_{i_t-\ell_j}$ and $i_t$ has
score $\beta_0$. Letters carry value $0$ and $i_t$ carries $\initscale^2A_V\hat e^{\mathrm g}_{s_{t-1}}$, so the head writes
$g(1-M_t)\hat e^{\mathrm g}_{s_{t-1}}$. Only $\beta$, $U$ and $W$ are trained.
\end{proof}

\begin{lemma}[Cyclic symmetry]
\label{lem:R-symmetry}
Throughout the schedule $RU=UR$ and $RW=WR$. Hence $w_k:=W_{s+k,s}$ does not depend on $s$, and
$\E_s\bigl[(\bm e_{s+1}-\bm e_s)^\top W\bm e_s\bigr]=w_1-w_0$.
\end{lemma}

\begin{proof}
Replacing $s_0$ by $s_0+1$ preserves the data distribution, shifts every count by one and leaves $M_t$ unchanged. Since the
cross-entropy satisfies $\ell(\bm z,y+1)=\ell(R^{-1}\bm z,y)$, the step loss is invariant under
$(U,W)\mapsto(R^{-1}UR,R^{-1}WR)$ for every $\beta$. This linear isometry therefore commutes with each stage's gradient field. Stage A
starts from $U=W=0$, which commutes with $R$, and each later stage starts where the previous one stopped, so uniqueness of solutions
keeps both blocks commuting with $R$.
\end{proof}

\subsection{The four-stage schedule}
Every stage is gradient flow $\tfrac{d\theta}{dt}=-P_k\nabla\Lstep(\theta)$, where $P_k$ projects onto the
trained block and all other parameters are frozen. The first three stages end at stopping times.
Write $\bm\pi_A(s)=\softmax(\initscale U\bm e_s)$ for the count-only prediction of the readout $U$.

\begin{center}
\small
\begin{tabular}{@{}llp{0.31\linewidth}p{0.35\linewidth}@{}}
\toprule
Stage & Train & Stop when & What it achieves \\
\midrule
A & $U$ & $\|\bm\pi_A(s)-\bm q_s\|_1\le\alpha_B$ for every $s$ & Fits the count-only prior \\[1mm]
B & $W$ & $\alpha=\alpha_B$ & Exposes the count--letter interaction \\[1mm]
C & $\beta$ & Leakage $\le\zeta$ on every prompt & Routes to the count and current letter \\[1mm]
D & $U,W$ & Continue with $\beta$ frozen & Gives positive margins and $\Lstep\to0$ \\
\bottomrule
\end{tabular}
\end{center}
Here $\alpha:=\operatorname{sign}(a)\,g\,(w_1-w_0)$ is the gate readout's coefficient on the move, and the readout and leakage scales are
\[
 \alpha_B=\frac{0.01\,p(1-p)|a|}{n^2},\qquad \zeta=\frac{0.01|a|}{\sqrt K\,B}.
\]
The logarithmic factor $B=O\bigl(K\log(nKd_{\mathrm h}/(p(1-p)\initscale))\bigr)$ is fixed in \cref{eq:R-Dconst}; it makes the final
readout robust to leakage. We write $\pi_0=\min(p,1-p)/K$ for the smallest probability of a local case $(s_{t-1},b_t)$.
Stage A starts at $U=0$, Stage B starts at $W=0$, and Stage C starts at $\beta=0$.
We require $|a|\le0.1$, $\initscale\le1$, and $\initscale\le\initscale_0$ as in \cref{ass:init}.

\label{lem:R-scale}
Every $0<\initscale\le\initscale_\star:=\min\{\initscale_0,1,(0.1\sqrt{d_{\mathrm h}})^{1/4}\}$ satisfies these
conditions, since $|a|=\initscale^4/\sqrt{d_{\mathrm h}}$. The structured initialization has $|\tilde b_m|=1$ on each of the $w$ MLP
units, so $\initscale_0=c\min(1,w^{-1/2})$ and $1/\initscale_\star\le C\sqrt w$.

\subsection{Local innovation identity}
Once Stage A has fit the count-only prior, the remaining gradient sees only how a feature responds to the current letter.

\begin{lemma}[Local innovation identity]
\label{lem:R-identity}
Write $s=s_{t-1}$. For a bounded $\Psi$ of $s$ and the letters, let $\bar\Psi=\E[\Psi\mid s,b_{-t}]$ and
$\partial_b\Psi=\Psi|_{b_t=1}-\Psi|_{b_t=0}$.
\begin{enumerate}[label=(\alph*),leftmargin=*,itemsep=3pt,topsep=2pt]
\item Among predictions $\bm\pi_s\in\Delta_K$ that depend only on $s$, the loss $\E[-\log\pi_{s}(s_t)]$ is minimized by $\bm\pi_s=\bm q_s$.
\item For every such $\bm\pi_s$,
\[
 \E\bigl[(\bm e_{s_t}-\bm\pi_s)\Psi\bigr]=\E\bigl[(\bm q_s-\bm\pi_s)\bar\Psi\bigr]+p(1-p)\,\E\bigl[(\bm e_{s+1}-\bm e_s)\,\partial_b\Psi\bigr].
\]
If $\Psi$ does not depend on $s$, then
$\E\bigl[(\bm e_{s_t}-\bm q_s)\Psi\bm x_t^\top\bigr]=\tfrac{p(1-p)}K\E[\partial_b\Psi]\,(R-I)$.
\item For probability vectors $\bm\pi,\bm\pi'$ and $\bm v\in\R^K$, $|(\bm\pi-\bm\pi')^\top\bm v|\le\tfrac12\|\bm\pi-\bm\pi'\|_1\operatorname{osc}(\bm v)$
and $\|\softmax(\bm z+\bm v)-\softmax(\bm z)\|_1\le\operatorname{osc}(\bm v)$.
\end{enumerate}
\end{lemma}

\begin{proof}
(a) Given $s$, the loss is the entropy of $\bm q_s$ plus $\KL(\bm q_s\,\|\,\bm\pi_s)$.

(b) Given $s$ and $b_{-t}$, $b_t$ is Bernoulli$(p)$ and $\Psi=\bar\Psi+(b_t-p)\,\partial_b\Psi$. Since $\E[\bm e_{s+b_t}\mid s]=\bm q_s$ and
$\bm\pi_s$ does not depend on $b_t$,
\[
 \E\bigl[(\bm e_{s+b_t}-\bm\pi_s)(b_t-p)\mid s\bigr]=p(1-p)(\bm e_{s+1}-\bm e_s),
\]
and averaging gives the identity. When $\Psi$ does not depend on $s$, the last term with $\bm x_t^\top$ appended averages over the
uniform $s$, independent of the letters, to $\tfrac{p(1-p)}K\E[\partial_b\Psi]\sum_s(\bm e_{s+1}-\bm e_s)\bm e_s^\top$, which is the stated matrix.

(c) Both vectors sum to one, so $\bm v$ may be shifted to have entries in $[-\tfrac12\operatorname{osc}(\bm v),\tfrac12\operatorname{osc}(\bm v)]$.
The softmax Jacobian $J=\operatorname{diag}(\bm\pi)-\bm\pi\bm\pi^\top$ gives $\|J\bm v\|_1=\sum_k\pi_k|v_k-\bm\pi^\top\bm v|\le\operatorname{osc}(\bm v)$;
integrate along $\bm z+\tau\bm v$, $\tau\in[0,1]$.
\end{proof}

At the prior, part (b) keeps only the response of a feature to the current letter, along the move $\bm e_{s+1}-\bm e_s$ of the rule
signal $\bm g_t$, with weight $p(1-p)=(1-\mixrate^2)/(2-2\cos\tfrac{2\pi}K)$; for $K=2$ this is $\tfrac14(1-\mixrate^2)$.

\begin{lemma}[Gate interaction derivatives]
\label{lem:R-derivatives}
Fix $\beta$ and $t$, and write $u=am_t$ for the current letter's interaction score, so $b_t=0$ gives $u=a$ and $b_t=1$ gives $u=-a$.
For every continuously differentiable $\Psi$ of $u$ and the other letters,
\begin{equation}
 \E[\partial_b\Psi]=-\E\Bigl[\int_{-a}^{a}\partial_u\Psi\,du\Bigr]=-\operatorname{sign}(a)\,2|a|\,\E_{\mathrm{aux}}[\partial_u\Psi].
 \label{eq:R-hat}
\end{equation}
Here the left side is averaged over $t$, and $\E_{\mathrm{aux}}$ averages over $u$ uniform on $[-|a|,|a|]$, over the other letters and over
$t$ uniform on $\{2,\ldots,n\}$. Write $\omega_0$, $\omega_{\mathrm{cur}}$ and $\omega_\delta$ for the gate head's attention weights on the count's own
position, the current letter and the letter at offset $\delta$, and $\omega_{\mathrm{oth}}=1-\omega_0-\omega_{\mathrm{cur}}$ for the total weight on all
other letters, so $M_t=\omega_{\mathrm{cur}}+\omega_{\mathrm{oth}}$. Then
\begin{equation}
 \begin{gathered}
 \frac{\partial M_t}{\partial\beta_0}=-\omega_0M_t,\qquad
 \frac{\partial M_t}{\partial\beta_n}=\omega_0\omega_{\mathrm{cur}},\qquad
 \frac{\partial M_t}{\partial\beta_\delta}=\omega_0\omega_\delta\quad(\delta\notin\{0,n\}),\\
 \frac{\partial M_t}{\partial u}=\omega_0\omega_{\mathrm{cur}},\qquad
 \frac{\partial\omega_{\mathrm{cur}}}{\partial u}=\omega_{\mathrm{cur}}(1-\omega_{\mathrm{cur}}),\qquad
 \frac{\partial\omega_\delta}{\partial u}=-\omega_\delta\,\omega_{\mathrm{cur}}\quad(\delta\ne n).
 \end{gathered}
 \label{eq:R-derivs}
\end{equation}
Every attention weight is a ratio of two sums whose terms change by a factor at most $e^{|a|}$ when $u$ or the other letters
change, so a product of $k$ weights changes by at most $e^{2k|a|}$ over $|u|\le|a|$; in particular
$\E_{\mathrm{aux}}[\omega_0\omega]\ge e^{-4|a|}\E[\omega_0\omega]$ for $\omega\in\{\omega_0,\omega_{\mathrm{cur}},\omega_\delta,\omega_{\mathrm{oth}}\}$, and $\E_{\mathrm{aux}}[\omega_0]\ge e^{-2|a|}\E[\omega_0]$.
\end{lemma}

\begin{proof}
The fundamental theorem of calculus between $u=a$ and $u=-a$ gives \cref{eq:R-hat}.
Differentiating the softmax ratios gives \cref{eq:R-derivs}; comparing their positive numerator and denominator terms gives the weight bounds.
\end{proof}

\subsection{Stage A: fit the count-only prior}
\begin{lemma}[Stage A]
\label{lem:R-A}
Let $\ell_p=\log(1/\min(p,1-p))$ and $L_A=\log(4K/\alpha_B^2)$. Stage A ends by
\[
 T_A\le\frac{2K\,(2\ell_p^2+KL_A^2)}{\initscale^2\alpha_B^2}.
\]
At its end, $\|\bm\pi_A(s)-\bm q_s\|_1\le\alpha_B$ for every $s$, and $\|\initscale U\|_F\le N_A:=3\sqrt K\,(2\ell_p^2+KL_A^2)^{1/2}$.
\end{lemma}

\begin{proof}
\emph{Reduction.} Here $W=0$, so $\bm z_t=\initscale U\bm x_t$ and each column $\bm v_s=\initscale U\bm e_s$ follows
\[
 \frac{d\bm v_s}{dt}=\frac{\initscale^2}K\bigl(\bm q_s-\softmax(\bm v_s)\bigr).
\]
This is gradient flow, at rate $\initscale^2/K$, on the convex $f_s(\bm v)=-\bm q_s^\top\bm v+\log\sum_ke^{v_k}$, whose excess over its infimum
is $\KL(\bm q_s\,\|\,\softmax(\bm v))$. For every comparator $\bar{\bm v}$, convexity gives
\[
 f_s(\bm v_s(t))\le f_s(\bar{\bm v})+\frac{K\|\bar{\bm v}\|^2}{2\initscale^2t},
 \qquad
 \|\bm v_s(t)-\bar{\bm v}\|^2\le\|\bar{\bm v}\|^2+\frac{2\initscale^2t}K\bigl(f_s(\bar{\bm v})-\inf f_s\bigr).
\]

\emph{Comparator.} Let $\bar{\bm v}$ have entries $\log(1-p)$ at $s$, $\log p$ at $s+1$ and $-L_A$ elsewhere. Then
$f_s(\bar{\bm v})-\inf f_s=\log(1+(K-2)e^{-L_A})\le\tfrac14\alpha_B^2$ and $\|\bar{\bm v}\|^2\le2\ell_p^2+KL_A^2$.

\emph{Time and norm.} At $t=T_A$ the first bound gives $\KL\le\tfrac12\alpha_B^2$, so Pinsker's inequality gives
$\|\softmax(\bm v_s)-\bm q_s\|_1\le\alpha_B$; the excess is nonincreasing along the flow, so the stopping time is at most $T_A$. At any
time $t\le T_A$ the second bound gives $\|\bm v_s(t)\|\le2\|\bar{\bm v}\|+(2\initscale^2T_A\alpha_B^2/(4K))^{1/2}\le3(2\ell_p^2+KL_A^2)^{1/2}$;
summing over the $K$ columns gives the bound on $\|\initscale U\|_F$.
\end{proof}

\subsection{Stage B: expose the interaction}
The goal is to give the gate readout a small coefficient on the move while keeping the routing fixed.
At zero bias, the innovation identity gives the following drive.

\begin{lemma}[Gate coefficient at zero bias]
\label{lem:R-coeff}
Let $n\ge2$, $|a|\le0.1$ and $\beta=0$. Then $\kappa:=p(1-p)\operatorname{sign}(a)\,\E[\partial_b(1-M_t)]\ge0.69\,p(1-p)|a|/n^2$.
\end{lemma}

\begin{proof}
By \cref{eq:R-hat}, $\operatorname{sign}(a)\E[\partial_b(1-M_t)]=2|a|\,\E_{\mathrm{aux}}[\partial_uM_t]$. For $|a|\le0.1$ the normalizer is
$Z=1+\sum_je^{am_j}\le1.61n$, so $\partial_uM_t=\omega_0\omega_{\mathrm{cur}}=e^u/Z^2\ge0.349/n^2$.
\end{proof}

\begin{lemma}[Stage B]
\label{lem:R-B}
During Stage B, $\max_k|gw_k|\le\alpha$ and
\[
 \frac{d\alpha}{dt}\ge0.92\,\frac{g^2p(1-p)|a|}{Kn^2}>0.
\]
Hence $\alpha$ reaches $\alpha_B$ by $T_B\le0.011\,Kg^{-2}$. From the end of Stage B until Stage D, for every $\beta$ and on every prompt,
\[
 \max_k|gw_k|\le\alpha_B,\qquad \operatorname{osc}(gW\bm e_s)\le2\alpha_B,\qquad \|\softmax(\bm z_t)-\bm\pi_A(s_{t-1})\|_1\le2\alpha_B .
\]
\end{lemma}

\begin{proof}
\emph{Drive.} Here $\beta=0$ and
\[
 \frac{dW}{dt}=g\,\E\bigl[(\bm e_{s_t}-\softmax(\bm z_t))(1-M_t)\bm x_t^\top\bigr].
\]
Split $\bm e_{s_t}-\softmax(\bm z_t)=(\bm e_{s_t}-\bm q_s)+(\bm q_s-\bm\pi_A(s))+(\bm\pi_A(s)-\softmax(\bm z_t))$. By \cref{lem:R-identity}(b) with
$\Psi=1-M_t$, the first part contributes $\tfrac gK\operatorname{sign}(a)\kappa(R-I)$. Column $s$ of the other two parts is $\tfrac gK$ times
a conditional mean of vectors with entries at most $\tfrac12\alpha_B$ and, by \cref{lem:R-identity}(c),
$\tfrac12\operatorname{osc}((1-M_t)gW\bm e_s)\le\max_k|gw_k|$. By \cref{lem:R-symmetry},
\[
 \frac{d(gw_k)}{dt}=\frac{g^2}K\bigl(\operatorname{sign}(a)\kappa\,([k=1]-[k=0])+r_k\bigr),
 \qquad |r_k|\le\bar r:=\tfrac12\alpha_B+\max_j|gw_j|.
\]

\emph{Growth with control of the other entries.} Thus $\tfrac{d\alpha}{dt}\ge\tfrac{g^2}K(2\kappa-2\bar r)$ and
$\bigl|\tfrac{d(gw_k)}{dt}\bigr|\le\tfrac{g^2}K(\kappa+\bar r)$. While $\max_j|gw_j|\le\alpha_B$, $\bar r\le1.5\alpha_B\le\kappa/3$ by
\cref{lem:R-coeff}, so $|\tfrac{d(gw_k)}{dt}|\le\tfrac43\tfrac{g^2}K\kappa\le\tfrac{d\alpha}{dt}$; since all start at $0$, $|gw_k|\le\alpha$. Up to the stopping
time $\alpha\le\alpha_B$, so a continuity argument keeps both bounds throughout Stage B, and
$\tfrac{d\alpha}{dt}\ge\tfrac43\tfrac{g^2}K\kappa\ge0.92\,g^2p(1-p)|a|/(Kn^2)$. Hence $T_B\le\alpha_BKn^2/(0.92g^2p(1-p)|a|)\le0.011Kg^{-2}$.

\emph{After Stage B.} $W$ is frozen until Stage D, so $\max_k|gw_k|\le\alpha_B$ and $\operatorname{osc}(gW\bm e_s)\le2\alpha_B$; the last bound is
\cref{lem:R-identity}(c) with $\bm v=(1-M_t)gW\bm x_t$, for every $\beta$.
\end{proof}

\subsection{Stage C: learn the routing}
Stage C must prove three things simultaneously: useful routing strength grows, irrelevant offsets are suppressed,
and the two useful sites remain balanced. Write
\[
 \bar\beta=\tfrac12(\beta_n+\beta_0)\ \ \text{(common routing strength)},
 \qquad
 \Delta=\beta_n-\beta_0\ \ \text{(routing imbalance)},
\]
and define the leakage by
\[
 \operatorname{Leak}(\beta)=\max_{t\ge2,\,\mathrm{prompt}}\omega_{\mathrm{oth}}.
\]
We work on the corridor $\mathcal C=\{\beta:|\Delta|\le1,\ \beta_\delta\le\bar\beta+0.1\text{ for }\delta\notin\{0,n\}\}$.

\begin{lemma}[The routing field]
\label{lem:R-field}
During Stage C, for every $\beta\in\mathcal C$:
\begin{enumerate}[label=(\alph*),leftmargin=*,itemsep=3pt,topsep=2pt]
\item \emph{Common strength grows and irrelevant biases fall:}
\[
 \frac{d\bar\beta}{dt}\ge0.42\,\frac{p(1-p)\alpha_B|a|}{n}\E[\omega_0\omega_{\mathrm{oth}}]>0,
 \qquad \frac{d\beta_\delta}{dt}\le0\quad(\delta\notin\{0,n\}).
\]
\item \emph{The useful sites remain balanced:} $\tfrac{d\Delta}{dt}<0$ at $\Delta=1$ and $\tfrac{d\Delta}{dt}>0$ at $\Delta=-1$.
\item \emph{Leakage controls the drive:}
\[
 \E[\omega_0\omega_{\mathrm{oth}}]\ge\frac{0.21}{n^2}\operatorname{Leak}(\beta),
 \qquad
 \frac{d\bar\beta}{dt}\ge0.0882\,\frac{p(1-p)\alpha_B|a|}{n^3}\operatorname{Leak}(\beta).
\]
\end{enumerate}
\end{lemma}

\begin{proof}
\emph{Preliminary mass bound.} On $\mathcal C$, $\beta_0,\beta_n\ge\bar\beta-\tfrac12$, $|\Delta|\le1$, and the other letters add at most
$(n-1)e^{|a|+\bar\beta+0.1}$ to the softmax normalizer $Z$. Hence
\[
 \frac{Z}{e^{\beta_0}}\le1+e^{1.1}+(n-1)e^{0.7}\le3.1n,
 \qquad\text{and likewise}\qquad
 \frac{Z}{e^{u+\beta_n}}\le3.1n .
\]

\needspace{14\baselineskip}
\emph{Bias dynamics.} For any relative offset $\delta$, \cref{lem:R-features} gives
$\tfrac{\partial\bm z_t}{\partial\beta_\delta}=-g\,\tfrac{\partial M_t}{\partial\beta_\delta}W\bm x_t$, so gradient flow gives
\[
 \frac{d\beta_\delta}{dt}=-\frac{\partial\mathcal L}{\partial\beta_\delta}
 =-\E\Bigl[\bigl(\bm e_{s_t}-\softmax(\bm z_t)\bigr)^\top gW\bm x_t\,\frac{\partial M_t}{\partial\beta_\delta}\Bigr].
\]
Split the residual as in Stage B. By \cref{lem:R-identity}(b) and \cref{lem:R-symmetry}, the innovation part equals
$-p(1-p)\,g(w_1-w_0)\,\E[\partial_b\partial_{\beta_\delta}M_t]$, and $g(w_1-w_0)=\operatorname{sign}(a)\alpha_B$. By \cref{lem:R-identity}(c) and
\cref{lem:R-A,lem:R-B}, on every prompt the prior part is at most $\tfrac12\alpha_B\cdot2\alpha_B$ and the readout part at most
$\tfrac12\cdot2\alpha_B\cdot2\alpha_B$, times $|\partial M_t/\partial\beta_\delta|$. With \cref{eq:R-hat},
\begin{equation}
 \frac{d\beta_\delta}{dt}=2p(1-p)\alpha_B|a|\,\E_{\mathrm{aux}}\Bigl[\frac{\partial^2M_t}{\partial u\,\partial\beta_\delta}\Bigr]+e_\delta,
 \qquad
 |e_\delta|\le3\alpha_B^2\,\E\Bigl[\Bigl|\frac{\partial M_t}{\partial\beta_\delta}\Bigr|\Bigr].
 \label{eq:R-bias}
\end{equation}
The velocity of each bias is a local interaction drive plus a small error. Both sides are linear in the bias coordinate, so
\cref{eq:R-bias} holds verbatim for $\bar\beta$ and $\Delta$, with $\tfrac{\partial M_t}{\partial\beta_\delta}$ replaced by the same
combination of derivatives. Since $n\ge2$, $\alpha_B=\tfrac{0.01\,p(1-p)|a|}{n^2}\le\tfrac{0.005\,p(1-p)|a|}{n}$.

\emph{Common strength.} Apply \cref{eq:R-bias} to $\bar\beta=\tfrac12(\beta_n+\beta_0)$. By \cref{eq:R-derivs},
$\tfrac12\bigl(\tfrac{\partial M_t}{\partial\beta_n}+\tfrac{\partial M_t}{\partial\beta_0}\bigr)=-\tfrac12\omega_0\omega_{\mathrm{oth}}$, whose $u$-derivative is
$\omega_0\omega_{\mathrm{oth}}\omega_{\mathrm{cur}}$, so with the preliminary mass bound
\begin{align*}
 \frac{d\bar\beta}{dt}&\ge2p(1-p)\alpha_B|a|\,\E_{\mathrm{aux}}[\omega_0\omega_{\mathrm{oth}}\omega_{\mathrm{cur}}]-\frac{3\alpha_B^2}{2}\,\E[\omega_0\omega_{\mathrm{oth}}]\\
 &\ge\alpha_B\,\E[\omega_0\omega_{\mathrm{oth}}]\Bigl(\frac{2e^{-0.4}}{3.1}-0.0075\Bigr)\frac{p(1-p)|a|}{n}.
\end{align*}
The numerical factor exceeds $0.42$.

\emph{Irrelevant offsets.} For $\delta\notin\{0,n\}$, \cref{eq:R-derivs} gives $\tfrac{\partial M_t}{\partial\beta_\delta}=\omega_0\omega_\delta$ with
$u$-derivative $-2\omega_0\omega_\delta \omega_{\mathrm{cur}}$, so \cref{eq:R-bias} and the preliminary mass bound give
\[
 \frac{d\beta_\delta}{dt}\le\alpha_B\,\E[\omega_0\omega_\delta]\Bigl(-\frac{4e^{-0.4}}{3.1}+0.015\Bigr)\frac{p(1-p)|a|}{n}\le0 .
\]

\emph{Imbalance.} Apply \cref{eq:R-bias} to $\Delta=\beta_n-\beta_0$: here $\tfrac{\partial M_t}{\partial\beta_n}-\tfrac{\partial M_t}{\partial\beta_0}=\omega_0(\omega_{\mathrm{cur}}+M_t)\le2\omega_0$, with $u$-derivative
$2\omega_0\omega_{\mathrm{cur}}(\omega_0-\omega_{\mathrm{cur}})$. At $\Delta=1$, $\tfrac{\omega_{\mathrm{cur}}}{\omega_0}=e^{u+1}\ge e^{0.9}$ for all
$|u|\le|a|$, so $\omega_0-\omega_{\mathrm{cur}}\le-0.59\,\omega_{\mathrm{cur}}$, and by the preliminary mass bound
\begin{align*}
 \frac{d\Delta}{dt}&\le-2p(1-p)\alpha_B|a|\,\E_{\mathrm{aux}}[1.18\,\omega_0\omega_{\mathrm{cur}}^2]+6\alpha_B^2\,\E[\omega_0]\\
 &\le\alpha_B\,\E[\omega_0]\Bigl(-\frac{2.36\,e^{-0.2}}{9.61}+0.06\Bigr)\frac{p(1-p)|a|}{n^2}<0.
\end{align*}
The case $\Delta=-1$ is symmetric, with $\tfrac{\omega_0}{\omega_{\mathrm{cur}}}\ge e^{0.9}$.

\emph{Leakage.} Let $\operatorname{Leak}(\beta)$ be attained at some position $t^*$. Every prompt at $t^*$ has $\omega_{\mathrm{oth}}\ge e^{-4|a|}\operatorname{Leak}(\beta)$, position $t^*$
has weight $\tfrac1{n-1}$, and $\omega_0\ge\tfrac1{3.1n}$; hence $\E[\omega_0\omega_{\mathrm{oth}}]\ge\tfrac{e^{-0.4}}{3.1\,n^2}\operatorname{Leak}(\beta)$.
The leakage factor exceeds $0.21$, so the combined drive has factor $0.0882$.
\end{proof}

\paragraph{Stage C hitting time.}
\begin{lemma}[Stage C]
\label{lem:R-C}
Throughout Stage C, $\beta\in\mathcal C$, $\bar\beta$ is nondecreasing, and $\beta_\delta\le0\le\bar\beta$ for every $\delta\notin\{0,n\}$. Stage C
ends by
\[
 T_C\le\frac{1134\,n^5\log(1.83n/\zeta)}{p^2(1-p)^2a^2\,\zeta},
\]
and at its end $|\Delta|\le1$ and, on every prompt and every $t\ge2$, $|M_t-\widetilde M_t|\le\zeta$ with
$\widetilde M_t=\sigma(am_t+\Delta)$.
\end{lemma}

\begin{proof}
\emph{Invariant corridor.} At $t=0$, $\beta=0$ lies in $\mathcal C$ with $|\Delta|<1$. Let $\tau$ be the first exit time from $\mathcal C$. On
$[0,\tau]$, \cref{lem:R-field}(a) gives $\bar\beta\ge0$ nondecreasing and $\beta_\delta\le0$, so $\beta_\delta-\bar\beta\le0<0.1$. An exit would
therefore have to occur through $|\Delta|=1$, which \cref{lem:R-field}(b) forbids, since there $\tfrac{d\Delta}{dt}$ points inward. Hence the
trajectory never leaves $\mathcal C$.

\emph{Hitting time.} On $\mathcal C$ with $\beta_\delta\le0$, $Z\ge e^{\beta_0}\ge e^{\bar\beta-1/2}$ gives
\[
 \operatorname{Leak}(\beta)\le(n-1)e^{0.1}e^{1/2}e^{-\bar\beta}\le1.83n\,e^{-\bar\beta},
\]
so Stage C ends once $\bar\beta\ge\log(1.83n/\zeta)$. Before then $\operatorname{Leak}(\beta)>\zeta$, and \cref{lem:R-field}(a,c) give
\[
 \frac{d\bar\beta}{dt}\ge0.0882\,\frac{p(1-p)\alpha_B|a|\zeta}{n^3}=0.000882\,\frac{p^2(1-p)^2a^2\zeta}{n^5},
\]
which gives the bound on $T_C$.

\emph{Approximation at exit.} Adding the other letters' weight $O$ to the numerator and denominator of $\widetilde M_t$ changes it by
at most $\tfrac OZ=\omega_{\mathrm{oth}}\le\zeta$.
\end{proof}

\subsection{Stage D: readout convergence after routing}
Once $|\Delta|\le1$ and leakage is at most $\zeta$, the gate output is a small perturbation of $(1-\widetilde M_t)\bm x_t$, which takes
two values per displayed count. A linear readout turns it into the transition, so the convex readout flow eventually gives
positive margins on every local case and drives the loss to zero. Measure readouts by $\|(U,W)\|_g:=|g|\,(\|U\|_F^2+\|W\|_F^2)^{1/2}$.

\begin{lemma}[Stage D]
\label{lem:R-D}
Write $\widetilde G(b)=1-\sigma\bigl(a(1-2b)+\Delta\bigr)$ for the routed own-count weight and $D=\widetilde G(1)-\widetilde G(0)$. Then
$\operatorname{sign}D=\operatorname{sign}a$ and $|D|\ge0.37|a|$, and the readout $(U^\star,W^\star)$ with
\[
 gW^\star=\frac{R-I}D,\qquad \initscale U^\star=I-\widetilde G(0)\,\frac{R-I}D
\]
has margin at least $\tfrac12$ on every prompt and every $t\ge2$, with $\|(U^\star,W^\star)\|_g\le\bar N:=7.84\sqrt K/|a|$. Measuring time from
the start of Stage D, there is $T_D\le\overline T_D$ (\cref{eq:R-Dconst}) such that every prediction has positive margin at every time
$t\ge T_D$, and $\Lstep\to0$.
\end{lemma}

\begin{proof}
\emph{Separator.} Since $|\Delta|\le1$ and $|a|\le0.1$, $D=\int_{-a}^{a}\sigma'(x-\Delta)\,dx$ has the sign of $a$ and $|D|\ge2|a|\sigma'(1.1)\ge0.37|a|$.
Write $G_t=1-M_t$ and $b=b_t$; since $\widetilde G(b)-\widetilde G(0)=bD$, \cref{lem:R-features} gives
\[
 \bm z^\star_t=\bm e_{s_{t-1}}+\frac{G_t-\widetilde G(0)}{D}\bigl(\bm e_{s_{t-1}+1}-\bm e_{s_{t-1}}\bigr)
 =\bm e_{s_t}+\frac{G_t-\widetilde G(b)}D\bigl(\bm e_{s_{t-1}+1}-\bm e_{s_{t-1}}\bigr).
\]
By \cref{lem:R-C}, $|G_t-\widetilde G(b)|\le\zeta$, so the margin is at least $1-2\zeta/|D|\ge\tfrac12$.
Also $\|gW^\star\|_F=\sqrt{2K}/|D|$ and $\|\initscale U^\star\|_F\le\sqrt K+\sqrt{2K}/|D|$, so with $\initscale\le1$,
$\|(U^\star,W^\star)\|_g\le\|\initscale U^\star\|_F+\|gW^\star\|_F\le\bar N$.

\emph{Local cases.} Write $\bar\gamma(s,b)$ for the margin of the logits $(\initscale U+\widetilde G(b)gW)\bm e_s$ at target $s+b$; it depends
only on the local case $(s_{t-1},b_t)=(s,b)$, since $\Delta$ does not depend on $t$. Each margin is $\bar\gamma(s,b)+\xi$ with
$|\xi|\le\zeta\operatorname{osc}(gW\bm e_s)\le2\zeta\|(U,W)\|_g$, and each case has probability at least $\pi_0$ at every $t$. If some margin is
nonpositive, then $\bar\gamma\le\max|\xi|$ for its case, so every prompt of that case has margin at most $2\max|\xi|$. Since each
cross-entropy term is at least $\log(1+e^{-\gamma_t})$,
\[
 \mathcal L\ge\pi_0\log\bigl(1+e^{-2\max|\xi|}\bigr).
\]

\emph{Convex-flow comparison.} The loss is convex in $v=(U,W)$, so for every comparator $\bar v$,
\[
 \mathcal L(v(t))\le\mathcal L(\bar v)+\frac{\|v(0)-\bar v\|^2}{2t},
 \qquad
 \|v(t)-\bar v\|\le\|v(0)-\bar v\|+\sqrt{2t\,\mathcal L(\bar v)},
\]
and, since every margin of $v^\star=(U^\star,W^\star)$ is at least $\tfrac12$, $\mathcal L(Cv^\star)\le(K-1)e^{-C/2}$. Thus a scaled separator makes
the loss arbitrarily small. We compare these bounds first over bounded times and then at long times.

\emph{Technical constants.} By \cref{lem:R-A,lem:R-B}, the readout at the start of Stage D has
$\|v(0)\|_g\le\initscale N_A+K\alpha_B\le\bar N_0:=1+\initscale N_A+K\alpha_B$. Set
\begin{equation}
 \begin{gathered}
 C_D=2\log\frac{10(K-1)}{\pi_0},\qquad
 \overline T_D=\frac{5(\bar N_0+C_D\bar N)^2}{\pi_0 g^2},\qquad
 R_D=\frac{4.3\sqrt{K-1}\,(\bar N_0^2+\bar N^2)}{\pi_0g^2},\\
 L_D=\log\bigl(100(K-1)R_D\bigr),\qquad
 T'=\bigl(100R_DL_D^2\bigr)^2,\qquad
 B=\max\bigl(C_D,\,2\log((K-1)(1+T')),\,\bar N_0+1.5\bigr).
 \end{gathered}
 \label{eq:R-Dconst}
\end{equation}
$\overline T_D$ bounds the burn-in, $T'$ is the time after which a second convexity argument certifies the margins, and $B$ sets the
leakage tolerance $\zeta=0.01|a|/(\sqrt KB)$ so that leakage cannot overturn a margin before $T'$. Note that $\bar N\ge1$, so
$\bar N_0+1.5\le B\le B\bar N$, and that $B\ge C_D\ge2\log40$ because $\pi_0\le\tfrac1{2K}$.

\emph{Initial loss bound.} With $\bar v=C_Dv^\star$, $\mathcal L(\bar v)\le0.1\pi_0$, so $\mathcal L(v(t))\le0.2\pi_0$ for
$t\ge T_D:=\tfrac{\|v(0)-C_Dv^\star\|^2}{0.2\pi_0}$, and $T_D\le\tfrac{(\bar N_0+C_D\bar N)^2}{0.2\pi_0g^2}=\overline T_D$.

\emph{Margins before $T'$.} For $0\le t\le T'$, take $\bar v=C_tv^\star$ with $C_t=\max(C_D,2\log((K-1)(1+t)))\le B$. Since
$t\,\mathcal L(\bar v)\le1$ and $|g|\le1$, the second convex-flow bound gives
\[
 \|v(t)\|_g\le\bar N_0+2B\bar N+1.5\le3B\bar N,
 \qquad\text{so}\qquad
 \max|\xi|\le6\zeta B\bar N\le0.48 .
\]
A nonpositive margin at a time $t\in[T_D,T']$ would then give $\mathcal L\ge\pi_0\log(1+e^{-0.96})\ge0.32\pi_0$, contradicting
$\mathcal L(v(t))\le0.2\pi_0$.

\emph{Margins after $T'$.} For $t\ge T'$, take $\bar v=Cv^\star$ with $C=2\log((K-1)(1+t))\ge1$, so $\mathcal L(\bar v)\le\tfrac1{1+t}$ and
$t\,\mathcal L(\bar v)\le1$. Then $\|v(t)\|_g\le B+2C\bar N$, and $2\max|\xi|\le4\zeta B+8\zeta C\bar N\le0.003+0.63\,C/B$. As
$\log(1+e^{-z})\ge e^{-z}\log2$ for $z\ge0$ and $1.26/B\le\tfrac12$, a nonpositive margin would give
\[
 \mathcal L(v(t))\ge\pi_0\log2\;e^{-0.003}\bigl((K-1)(1+t)\bigr)^{-1.26/B}\ge0.69\,\pi_0(K-1)^{-1/2}(1+t)^{-1/2}.
\]
The first convex-flow bound with the same $\bar v$ gives instead
\[
 \mathcal L(v(t))\le\frac1{1+t}+\frac{(\bar N_0+C\bar N)^2}{2g^2t}\le\frac1{1+t}+\frac{C^2(\bar N_0^2+\bar N^2)}{g^2t}.
\]
Since $T'\ge R_D^2\ge(4.3\sqrt{K-1}/\pi_0)^2$, the first term is at most $0.33\,\pi_0(K-1)^{-1/2}(1+t)^{-1/2}$, and the second is at
most $0.33\,\pi_0(K-1)^{-1/2}(1+t)^{-1/2}$ once $C^2/\sqrt t\le1/R_D$. With $\tau=\sqrt t$, $C\le2\log(2(K-1))+4\log\tau$. At
$\tau'=\sqrt{T'}=100R_DL_D^2$ one has $\log(2(K-1))\le L_D$ and $\log\tau'\le\log(100R_D)+2\log L_D\le2L_D$, so
$C^2/\tau'\le(10L_D)^2/(100R_DL_D^2)=1/R_D$; the left side decreases in $\tau$ beyond $\tau'$, so the condition holds for all $t\ge T'$.
The two bounds on $\mathcal L(v(t))$ contradict each other.

\emph{Convergence.} Every margin is therefore positive at every $t\ge T_D$, and $\limsup_t\mathcal L(v(t))\le(K-1)e^{-C/2}$ for every $C$,
so $\Lstep\to0$.
\end{proof}

\subsection{Proof of the theorem: assembly}
\begin{proof}[Proof of \cref{thm:routing}]
Let $T=T_A+T_B+T_C+T_D$.

\emph{Entry.} \Cref{lem:R-A,lem:R-B,lem:R-C} bound the first three stopping times. At the end of Stage C,
\cref{lem:R-C} gives
\[
 |\Delta|\le1,\qquad |M_t-\widetilde M_t|\le\zeta,
\]
so \cref{lem:R-D} gives positive margins at every time $t\ge T$.

\emph{Persistence.} Stage D freezes $\beta$, so the leakage stays at most $\zeta$ and $|\Delta|\le1$; the trajectory is in
$\mathcal G_\zeta$ at every $t\ge T$.

\emph{Convergence.} Convex readout flow gives $\Lstep\to0$ by \cref{lem:R-D}.

\emph{Time.} With $|a|=\initscale^4/\sqrt{d_{\mathrm h}}$, $|g|=\initscale^2$ and $\alpha_B=0.01\,p(1-p)|a|/n^2$, the stage bounds are
\[
 \begin{aligned}
 T_A&=O\Bigl(\frac{K^2n^4d_{\mathrm h}L_A^2}{p^2(1-p)^2\initscale^{10}}\Bigr), &
 T_B&=O(K\initscale^{-4}),\\
 T_C&=O\Bigl(\frac{\sqrt K\,n^5B\log(nKB/|a|)}{p^2(1-p)^2|a|^3}\Bigr), &
 T_D&\le\overline T_D=O\Bigl(\frac{K^2d_{\mathrm h}(\bar N_0+C_D)^2}{\min(p,1-p)\,\initscale^{12}}\Bigr).
 \end{aligned}
\]
Here $L_A$, $\bar N_0$, $C_D$ and $B$ are $O\bigl(K\log(nKd_{\mathrm h}/(p(1-p)\initscale))\bigr)$, so the sum is
$T\le\mathrm{poly}(n,K,d_{\mathrm h},\initscale^{-1},(p(1-p))^{-1})$.
\end{proof}

\subsection{Reference solution and controls}
The positive proof above learns the useful routing. We now give a reference executor and show why the seed direction matters,
then compare outcome training from the same initialization.

\paragraph{Reference gate executor.}
\begin{proof}[Proof of \cref{prop:gate}]
\emph{The routing.} Set
\[
 \bar\beta^\star=\log\frac{10n}{|a|},\qquad \beta^\star_0=\beta^\star_n=\bar\beta^\star,\qquad \beta^\star_\delta=0\ \text{otherwise}.
\]
By \cref{lem:R-features}, with $O\le(n-1)e^{|a|}$ the other letters' weight, the attention outside the two sites is
\[
 \frac{O}{e^{\bar\beta^\star}+e^{am_t+\bar\beta^\star}+O}\le Oe^{-\bar\beta^\star}\le0.111|a|,
\]
and it changes $M_t$ from $\widetilde M_t=\sigma(am_t)$ by at most that much.

\emph{The readout.} Take $(U^\star,W^\star)$ of \cref{lem:R-D} with $\Delta=0$, for which $D=\tanh(a/2)$ and $|D|\ge0.499|a|$. As in its
proof, the leakage changes every margin by at most $2\cdot0.111|a|/|D|<\tfrac12$, so every margin is at least $\tfrac12$ and
$\theta^\star_{\mathrm{gate}}\in\mathcal G_{0.111|a|}$. The stages of \cref{thm:routing} reach $\mathcal G_\zeta$ with $\zeta=0.01|a|/(\sqrt KB)$ by
\cref{lem:R-C,lem:R-D}.
\end{proof}

\paragraph{Wrong-seed control.}
\begin{lemma}[Wrong seed]
\label{lem:wrong-seed}
Replace the key seed $W_K\hat e_\ell=\initscale A_K\hat e_{\mathrm h}$ by $W_K\hat e_{0}=\initscale A_K\hat e_{\mathrm h}$, so that the key reads the direction
both letters share. Then, for every $\beta$ and every readout:
\begin{enumerate}[label=(\roman*),leftmargin=*,itemsep=2pt,topsep=2pt]
\item $\bm z_t$ is a function of $s_{t-1}$ and $t$ alone;
\item $\Lstep\ge-p\log p-(1-p)\log(1-p)$, and for each displayed count one of the two letter values has nonpositive margin;
\item $\partial_b(1-M_t)=0$, so at the prior the gate-readout gradient of \cref{lem:R-identity}(b) vanishes.
\end{enumerate}
The staged schedule, and any training of $\beta$ and the readout, therefore stays at or above the count-only prior.
\end{lemma}

\begin{proof}
\emph{Representation restriction.} Every letter now has key $\initscale^2A_K\hat e_{\mathrm h}$, so the scores at $i_t$ are
$a+\beta_{i_t-\ell_j}$ and $M_t$ depends on $t$ and $\beta$ only. By \cref{lem:R-features}, this gives (i).

\emph{Loss consequence.} Given $s_{t-1}=s$, the target is $s$ or $s+1$ with probabilities $1-p$ and $p$, so by \cref{lem:R-identity}(a)
the loss is at least the entropy of $\bm q_s$. The two targets have margins at most $z_s-z_{s+1}$ and $z_{s+1}-z_s$, one of which is
nonpositive. This gives (ii).

\emph{Zero interaction coefficient.} Since $M_t$ does not depend on $b_t$, $\partial_b(1-M_t)=0$, and \cref{lem:R-identity}(b) gives (iii).
\end{proof}

\paragraph{Outcome training in the same model.}

\begin{corollary}[Outcome training from the structured initialization]
\label{cor:routing-outcome}
Take the architecture of \cref{thm:routing}, with any fixed mask for the gate head at the answer position, and let
$\mixrate$ be as in \cref{eq:lambda}. From the structured initialization, or from it plus any perturbation with zero readout under which
\cref{ass:init} still holds, gradient flow on $\Lout$ with $\beta$ fixed obeys the plateau bound of \cref{thm:outcome-end}
for $\initscale\le\min(\initscale_\star,c\mixrate)$ (and \cref{rem:lambda-zero} when $\mixrate=0$).
\end{corollary}

\begin{proof}
The structured initialization satisfies \cref{ass:init} with $a_0=\sqrt K$, $\varsigma=1$ and $|\tilde b_m|=1$, and so does any
admissible perturbation; the readout is zero. A fixed mask is allowed by \cref{ass:arch}: each attention row is still a
softmax over its visible sites, zero-score attention is still uniform over them, and a mask can only lower the numbers of
scores and sites, so the arguments of \cref{prop:score-degree}, \cref{prop:network}(i)--(iii), \cref{lem:score-sum,lem:end-smooth}
and \cref{thm:outcome-end} apply unchanged.
\end{proof}

\begin{remark}[Scope]
\label{rem:routing-scope}
The step guarantee trains $\beta$ in a stage of its own, while the plateau of \cref{cor:routing-outcome} holds $\beta$ at zero.
At $K=2$, $\initscale=0.02$, $p=\tfrac14$ ($\mixrate=\tfrac12$) and $d_{\mathrm h}=1$ the bound on $T$, dominated by $T_C$, exceeds
$\mixrate^{-n}$ unless $n>138$; with $\initscale=\min(\initscale_\star,c\mixrate)$, which does not depend on $n$, the separation is
$\mathrm{poly}(n,d_{\mathrm h},w)$ against $c\mixrate^{-n}$ for each fixed $K$. The result covers transitions $t\ge2$, since the first
prediction has no count at its own position.
\end{remark}
